%!TEX root = main.tex

A description of the formalization of InferTypes, InferWidths, InferResets, ExpandConnects, ExpandWhens, LowerTypes, and RemoveResets. 



\subsubsection{``Infer Widths''}
Implicit widths are allowed to be defined in HiFIRRTL programs.
Such unspecified widths are infered by the Infer Widths pass. 
The general rules are,
\begin{itemize}
\item The only allowed width for the explicit defined width is itself.
\item The implicit width will be assigned by the largest width that can
  be found through signal connections.
\item The last connection semantics is not considered at this stage,
  all assigned signals are taking into account.
\end{itemize}

Following the rules, we first introduce an additional map $wm$, which
maps from variables to types. Then we scan the statements in $ss$ to
find components with implicit width and create the entries for
them in $wm$. For one statement which declares a component $v$,
the {\tt infer\_width} function will update the $wm$ by
  \begin{gather}
    \inferii[]
    {\seqq{ce \vdash {\tt infer\_width} (declare\;v, wm) \rightsquigarrow wm'}}
    {\seqq{ce \vdash {\tt width}(t_v) = 0}} {\seqq{wm':= wm \uplus {v : t_v}}}.
    \end{gather}
The width is inferred according to connection statements.
Since the last connection semantics is not considered, each
connection to the element provides the potential width. A function
named {\tt max\_width} compares the potential widths and outputs
the maximum one. For one statement which connects an expression $e$
to a component $v$, the {\tt infer\_width} function will update $wm$ by
  \begin{gather}
    \inferiii[]
    {\seqq{ce \vdash {\tt infer\_width }(v<=e, wm) \rightsquigarrow wm'}}
    {\seqq{wm \vdash v : t_v}}
    {\seqq{ce \vdash t_e:= {\tt type\_of\_expr }(e)}}
    {\seqq{wm':= wm \uplus {v : {\tt max\_width }(t_v,t_e)}}}.
    \end{gather}
After applying {\tt infer\_width} on all the statements in $ss$,
the $wm$ contains the inferred width for every implicit defined
component. By using {\tt map2} function merges the $ce$ and $wm$
with the same entry, the component environment $ce$ now contains the
result of the inferred width.

Instead of applying {\tt infer\_width} directly on $ce$, the map
$wm$ contains only type information of the implicit ones that can be
manipulated more efficiently.

%% \begin{description}
%% \item[The input] of InferWidth takes the statments and the component
%%   environment $ce$ that passed from the InferType.  In $ce$ the
%%   unspecified widths are recorded by zero width .  The statements
%%   are assumed to contain only legal connections, which means the two
%%   sides of the connection have equivalent type (mentioned?) and the
%%   flow is from a smaller width to a lager one.
%% \item[The output] of InferWidth give a new component environment $ce'$
%%   where all unspecified types have been infered.  Still it is possible
%%   that after this pass, signals are infered to zero width.  The overall
%%   compilation will use other passes to deal with zero width, which is
%%   not concerned here.
%% \end{description}

%% To state the correctness of InferWidth that implement in Coq, we first
%% recall the semantics of width infering.  There are statements that
%% declare new component, or the statements that define connections. For
%% those components declared with explicit widths, we state that the
%% widths will be kept unchange after the pass.  The implicit width
%% declarations will be collected and await to be assigned to a proper width.
%% The connections give hints to the infering process. For example,
%% \begin{lstlisting}
%%   wire c: SInt
%%   c <= shl(SInt("b-100101100101"), 1)
%% \end{lstlisting}
%% the width of wire $c$ is implicit defined. The connection from an
%% shift-left expression to $c$ will provide the protential width for $c$.
%% According to the rules, the width of c now is 13, because the
%% shift-left expression is of type \lstinline!SInt<13>!. Before applying last
%% connect semantics, the module may contain more than one connections to
%% one target. For example,
%% \begin{lstlisting}
%%   wire c: SInt
%%   wire d: SInt<14>
%%   c <= shl(SInt("b-100101100101"), 1)
%%   c <= d
%% \end{lstlisting}
%% Now, wire $c$ has the width as same as $d$, which is 14. The latter connection
%% from d gives a larger width to the destination $c$. 

\subsubsection{``Expand Connects''}
In HiFIRRTL, connections can be made between aggregate types.  These
connections can be either full or partial.  The purpose of this pass
is to break down connections between aggregate types and produce
individual ones of corresponding ground types. The general rules are,
\begin{itemize}
\item Connections between ground types can not be expanded.
\item Connections between aggregate types will be expanded recursively.
\item The direction of the connection will be flipped,
      if the expanded field in a bundle type is declared \prettyll{flip}
\end{itemize}

The pass focuses on connection statements in $ss$. The left-hand side of a
connection is a reference to a component and the right-hand side is an
expression. 
%% For example,
%% \begin{lstlisting}
%% module MyModule :
%%   input myinput: {a: UInt<1> [1]}
%%   output myoutput: {a: UInt<1> [2], b: UInt<1>}
%%   myoutput <- myinput
%% \end{lstlisting}
%% the connection between myinput and myoutput is a partial connection of
%% aggregate types. It can then be expand through ExpandConnects pass and get,
%% \begin{lstlisting}
%% module MyModule :
%%   input myinput_a: UInt<1>
%%   output myoutput_a_0: UInt<1>
%%   output myoutput_a_1: UInt<1>
%%   output myoutput_b: UInt<1>
%%   myoutput_a_0 <= myinput_a_0
%% \end{lstlisting}

%% \begin{description}
%% \item[The input] of ExpandConnects consists of the sequence of
%%   statments representing the original program, the result component
%%   enviroment $ce3$ from InferWidths pass. The $ce3$ collects the
%%   component identifies to their types.
%% \item[The output] of ExpandConnects contains a program in which all
%%   connections between aggregate types are replaced by connections of
%%   their sub elements and a corresponding $ce4$. In $ce4$, new
%%   identifiers added and their types are consistent to the flattened
%%   aggregate type. As the aboved example, ouput port \lstinline!myport!
%%   is replaced by 3 output ports of ground types.
%% \end{description}

\subsubsection{``Expand Whens''}

HiFIRRTL allows to define conditional components, mostly connect
statements that depend on a condition.  However, a circuit cannot
contain conditional connections; they need to be replaced by
multiplexers.  The ``expand whens'' transformation does exactly this:
it replaces connects within \textit{when} statements by multiplexers
where necessary.  For example, \pagebreak[1]\\
\begin{minipage}[t]{12.49em}
\vspace*{-0.94\baselineskip}
\begin{lstlisting}
module TestCondition
  input in: UInt<32>
  input condition: UInt<1>
  output posv: UInt<32>
  output negv: SInt<33>
  posv <= 0
  when condition:
    posv <= in
  else:
    wire a: SInt<33>
    a <= neg in
    negv <= a
\end{lstlisting}
\end{minipage}
is translated to\qquad{}%
\begin{minipage}[t]{16.81em}
\vspace*{-0.94\baselineskip}
\begin{lstlisting}
module TestCondition
  input in: UInt<32>
  input condition: UInt<1>
  output posv: UInt<32>
  output negv: SInt<33>
  wire a: SInt<33>
  posv <= mux(condition, in, 0)
  a <= neg in
  negv <= validif(not condition, a)
\end{lstlisting}
\end{minipage}

Additionally, ``expand when'' removes repeated connects,
thereby implementing the \emph{last connect} semantics.
The ``expand when'' transformation produces a map
where every connected component is mapped to the one expression
that it is connected to.
This map can then easily changed back into a list of connect statements,
as shown in the example above.
Note that this phase only touches on (full) connect statements;
node statements are left unchanged.

The specification of the Expand Whens transformation is therefore:
\begin{description}
\item[The input] of Expand Whens is the code of a module
where all types of components have been lowered to ground types.
For example, a statement \lstinline!reg a: { a1: UInt<32>, a2: SInt<16> }, clk!
is already translated to the two statements
\lstinline!reg a$a1: UInt<32>, clk! and
\lstinline!reg a$a2: SInt<16>, clk!.
Also, all partial connects (\lstinline!<-!) have been replaced by appropriate full connects (\lstinline!<=!).

\item[The output] of Expand Whens is a pair (list of statements, map of connects);
the list of statements does not contain any connect statements
(but mainly declarations of registers and wires and definitions of nodes,
copied verbatim from the input),
while the map of connects contains
the last connect found on each path through the input statement sequence that declares the component;
if multiple execution paths declare the component but provide different definitions,
then a suitable multiplexer or \lstinline!validif! expression is returned.
\end{description}

In the above example, the resulting code always connects \lstinline!a <= in!
because wire \lstinline!a! is only declared in the \lstinline!else! path.
On the other hand, output \lstinline!negv! is connected using a \lstinline!validif! statement
because an output is declared in every path.
(In a real circuit, this is probably a designer error,
and a FIRRTL compiler is required to produce an error message.
The error message can be suppressed by a statement \lstinline!negv is invalid!.
Our formalization of the lowering transformation does not produce this error message
but leaves it to the LoFIRRTL compiler.)
The output \lstinline!posv! is connected twice;
this is translated to a single connect with a multiplexer.

In the Coq implementation, the result of ``\texttt{find v}''
(try to find the definition of \texttt{v} in a map),
when applied to the resulting map,
produces \texttt{None} if the component has not been declared;
\texttt{Some R\_default} if the component has been declared but not defined
(i.e.\@, not connected to);
and \texttt{Some (R\_fexpr e)} if the last connect assigned expression \texttt{e} to \texttt{v}.
At the end of a \lstinline!when! statement,
the intermediary maps of the two execution paths are joined into one.
If the component has been declared in both paths%
---this means that the component has been declared before the \lstinline~when~ statement%
---but only defined in one, a \lstinline!validif! expression is produced.
If the component is defined in both paths
and the definitions differ, a \lstinline!mux! expression is produced.

